% Part: normal-modal-logic
% Chapter: completeness
% Section: introduction

\documentclass[../../../include/open-logic-section]{subfiles}

\begin{document}

\olfileid[pt]{nml}{com}{int}

\olsection{Introdução}

Se $\Sigma$ é um sistema modal, então o teorema da correção estabelece
que, se $\Sigma \Proves !A$, $!A$ é válida em qualquer classe
$\mClass{C}$ de modelos na qual todas as instâncias de todas as
!!{formula}s em~$\Sigma$ são válidas. Em particular, isso significa
que, se $\Log{K} \Proves !A$, então $!A$ é verdadeira em todos os
modelos; se $\Log{KT} \Proves !A$, então $!A$ é verdadeira em todos os
modelos reflexivos; se $\Log{KD} \Proves !A$, então $!A$ é verdadeira
em todos os modelos seriais; e assim por diante.

A completude é a recíproca da correção: dizer que \Log{K} é completa
significa, por exemplo, que, se !!a{formula}~$!A$ é válida, então
$\Proves !A$. Demonstrar a completude é muito mais difícil do que
demonstrar a correção. Para começar, é útil considerar a
contrapositiva: \Log{K} é completa se e somente se, sempre que
$\Proves/ !A$, existe um contramodelo, isto é, um modelo~$\mModel{M}$
tal que $\mSat/{M}{!A}$. Equivalentemente (negando $!A$), poderíamos
demonstrar que, sempre que $\Proves/ \lnot !A$, existe um modelo
de~$!A$. Na construção de tal modelo, podemos usar as informações
contidas em $!A$. Quando encontramos modelos para !!{formula}s
específicas, muitas vezes fazemos o mesmo: por exemplo, se queremos
encontrar um contramodelo de $p \lif \Box q$, sabemos que ele deve
conter um mundo no qual $p$ é verdadeiro e $\Box q$ é falso. E um
mundo no qual $\Box q$ é falso exige que haja um mundo acessível a
partir dele no qual $q$ é falso. Isso é tudo o que precisamos saber:
quais mundos tornam verdadeiras quais !!{propositional variable}s e
quais mundos são acessíveis a partir de quais outros mundos.

No caso da demonstração da completude, contudo, não temos uma
!!{formula}~$!A$ específica para a qual estamos construindo um modelo.
Queremos estabelecer que existe um modelo para todo $!A$ tal que
$\Proves/[\Sigma] \lnot !A$. Esse é um requisito mínimo, pois,
\emph{se} $\Proves[\Sigma] \lnot !A$, pela correção não existe modelo
de~$!A$ (no qual $\Sigma$ seja verdadeiro). Observe agora que
$\Proves/[\Sigma] \lnot !A$ se e somente se $!A$ é
$\Sigma$-consistente. (Lembre-se de que $\Sigma
\Proves/[\Sigma] \lnot !A$ e $!A \Proves/[\Sigma] \lfalse$ são
equivalentes.) Portanto, nossa tarefa é construir um modelo para toda
!!{formula} $\Sigma$-consistente.

O recurso que usaremos consiste em encontrar um conjunto
$\Sigma$-consistente de !!{formula}s que contenha~$!A$, mas também
outras fórmulas que nos digam como deve ser o mundo que torna $!A$
verdadeira. Esses conjuntos são chamados \emph{conjuntos
$\Sigma$-consistentes completos}. Não basta construir um modelo com um único mundo que
torne~$!A$ verdadeira: ele terá de conter vários mundos e uma relação
de acessibilidade. O conjunto $\Sigma$-consistente completo que
contém~$!A$ também conterá outras !!{formula}s da forma $\Box !B$
e~$\Diamond !C$. Em todos os mundos acessíveis, $!B$ deve ser
verdadeira; em pelo menos um deles, $!C$ deve ser verdadeira. Para
conseguir isso, tomaremos simplesmente \emph{todos} os conjuntos
$\Sigma$-consistentes completos possíveis como base do conjunto de
mundos. Uma parte delicada será determinar quando um conjunto
$\Sigma$-consistente completo deve contar como acessível a partir de
outro em nosso modelo.

Mostraremos que, no modelo assim definido, $!A$ é verdadeira em um
mundo---que também é um conjunto $\Sigma$-consistente completo---se e
somente se $!A$ é !!a{element} desse conjunto. Se $!A$ é
$\Sigma$-consistente, ela será !!a{element} de pelo menos um conjunto
$\Sigma$-consistente completo (um fato que demonstraremos) e, assim,
haverá um mundo no qual $!A$ é verdadeira. Portanto, teremos um único
modelo no qual toda !!{formula} $\Sigma$-consistente~$!A$ é verdadeira
em algum mundo. Esse único modelo é o modelo \emph{canônico}
de~$\Sigma$.

\end{document}
